\documentclass[11pt]{article}
\usepackage[a4paper,margin=23mm]{geometry}
\usepackage{amsmath,amssymb,amsthm,booktabs,microtype,xurl,hyperref}
\hypersetup{colorlinks=true,linkcolor=blue,urlcolor=blue,citecolor=blue,
pdftitle={P54 coordinate choices and threshold-aware convergence}}
\setlength{\emergencystretch}{2em}
\newcommand{\diag}{\operatorname{diag}}
\newcommand{\cond}{\operatorname{cond}}
\newtheorem{proposition}{Proposition}
\title{Route-F P54: What a New Coordinate Can Improve\\
Units, Metric Whitening and Threshold-Aware Series}
\author{Frozen-action convergence audit}
\date{17 September 2026}
\begin{document}
\maketitle
\begin{abstract}
We test the suggestion of absorbing $\pi$ into a coordinate unit, and
distinguish linear-solver convergence, momentum-series convergence and
quantum loop control. Coordinate changes leave the relative loop/tree
spectrum invariant. Metric whitening improves the conditioning of the
known local kinetic matrix but does not repair its large insertion.
A threshold-derived conformal variable genuinely extends convergence of
an equal-mass one-loop master. We derive its coefficients and a rigorous
tail bound, then test actual masses of the unchanged P54 benchmark.
This is a new controlled representation of a known kernel, not a new
physical action or an all-orders convergence proof.
\end{abstract}

\section{Putting pi into the ruler}
The proposed coordinate is $u=x/\pi$, so that $x=\pi$ becomes $u=1$.
If $x$ is a dimensional length, a dimensionless coordinate instead reads
$u=x/(\pi\ell_0)$ for a specified reference length $\ell_0$.
Neither operation changes a dimensionless physical ratio.
For example, on a circle of radius $R$, let $u=\theta/\pi$. Then
\begin{equation}
 u\sim u+2,\quad ds^2=\pi^2R^2du^2,\quad
 -\Delta=-\frac1{\pi^2R^2}\partial_u^2,\quad
 -\Delta e^{in\pi u}=\frac{n^2}{R^2}e^{in\pi u}.
\end{equation}
The explicit $\pi$ moves into the metric and Fourier modes; the spectrum
does not move. Treating a noncompact radial field as a periodic angle,
in contrast, would change the configuration space and therefore the theory.

The same conclusion holds for a spacetime unit change. In four Euclidean
dimensions take $x=\pi u$, $\psi(u)=\pi\phi(\pi u)$. Then
\begin{align}
 S&=\int d^4x\left[\tfrac12(\partial_x\phi)^2+
               \tfrac12m^2\phi^2+\frac{\lambda}{4!}\phi^4\right]\\
  &=\int d^4u\left[\tfrac12(\partial_u\psi)^2+
               \tfrac12(\pi m)^2\psi^2+\frac{\lambda}{4!}\psi^4\right].
\end{align}
Momentum and subtraction scale become $p_u=\pi p_x$, $\mu_u=\pi\mu_x$.
Thus $p^2/m^2$, $\log(m^2/\mu^2)$ and the dimensionless quartic are
unchanged. A convention $\widehat\lambda=\lambda/(16\pi^2)$ similarly
relocates loop factors without removing them from predictions.

\section{The invariant test for field-coordinate changes}
For radial fluctuations use the existing tree metric and known insertion
\begin{equation}
 G=\diag(12/5,2,1),\qquad Z=\Delta Z_{S,\mathrm{hard}}+
                                      \Delta Z_{\mathrm{gauge}},
\end{equation}
from P-MOM1. This $Z$ excludes the scalar soft--soft nonlocal term and
unfitted fermions; it is not a physical residue. Its generalized spectrum is
\begin{equation}
 \operatorname{eig}(Z,G)=(.0089951842,.2733736138,2.4665928386).
 \label{eq:relative}
\end{equation}
For an invertible constant field map $\delta r=Jq$,
\begin{equation}
 G'=J^TGJ,\quad Z'=J^TZJ,\quad K'(s)=J^TK(s)J.
\end{equation}
Consequently
\begin{align}
 \det(Z'-\lambda G')&=(\det J)^2\det(Z-\lambda G),\\
 (G')^{-1}Z'&=J^{-1}(G^{-1}Z)J.
\end{align}
Every generalized eigenvalue in \eqref{eq:relative}, and hence
$\rho=\max|\lambda|$, is invariant. The same determinant identity
applies to $K(s)$. For a source term $j^T\delta r$, $j'=J^Tj$ and
\begin{equation}
 j'^T(K')^{-1}j'=j^TK^{-1}j.
\end{equation}
These are exact finite-matrix identities, not numerical approximations.
For $J=\pi I$, both tree and loop matrices acquire $\pi^2$.
For $J=I/\pi$, both get smaller by $\pi^2$. Comparing the loop matrix
with the \emph{old} tree metric after either transformation is inconsistent.
General quantum field redefinitions must also transform interactions,
sources, matching coefficients and the regulated measure consistently;
higher-order truncations require particular care \cite{Criado}.

\subsection{A nonlinear chart cannot evade the stationary Hessian}
For $r=f(q)$, the ordinary potential Hessian obeys
\begin{equation}
 \partial_a\partial_b V(f(q))
 =J^i{}_a H_{ij}J^j{}_b+t_i\,\partial_a\partial_b f^i,\qquad
 t_i=\partial_iV.
\end{equation}
The extra term disappears at a stationary point. Off shell one must
include it, or use the covariant Hessian associated with the transformed
kinetic metric; the connection term removes it.
For the local positive-radius chart $r_i=r_{0i}e^{q_i}$,
\begin{equation}
 J=\diag(r_i),\quad
 H^{\mathrm{ordinary}}_{q}=J^TH_rJ+\diag(r_it_i),\quad
 \Gamma^i{}_{ii}=1,\quad
 \nabla_q^2V=J^TH_rJ.
\end{equation}
The existing loop tadpoles and common finite mass counterterms verify
these identities and $t+H_{\rm CT}r_0=0$ numerically. A coordinate
curvature obtained by omitting the off-shell gradient term is not a
new stability result. Sylvester's inertia law prevents an invertible
congruence from turning a stationary negative mode into a positive one.

\section{A real numerical benefit: metric preconditioning}
The known local matrix $M=G+Z$ is positive definite, even though the
full momentum-dependent background kernel can be indefinite.
For this \emph{local matrix only}, compare
\begin{equation}
 J_{\rm tree}=G^{-1/2},\qquad J_{\rm local}=M^{-1/2}.
\end{equation}
The first makes the tree metric the identity. The second makes the
known local kinetic matrix the identity, but does not make the
relative loop/tree eigenvalues smaller.

\begin{center}
\begin{tabular}{lrr}
\toprule
Chart & $\cond_2(J^TMJ)$ & $\rho(Z',G')$\\
\midrule
Original & 6.907906 & 2.466593\\
$u=x/\pi$ & 6.907906 & 2.466593\\
$u=\pi x$ & 6.907906 & 2.466593\\
Relative-VEV / local log chart & 50.581906 & 2.466593\\
Tree-metric whitening & 3.435688 & 2.466593\\
Known-local-metric whitening & 1 & 2.466593\\
\bottomrule
\end{tabular}
\end{center}
Thus setting all background coordinates to order one is not
automatically good conditioning. No coupling or background has changed.

For an SPD linear system $Ax=b$, Richardson iteration with
$\alpha=2/(\lambda_{\min}+\lambda_{\max})$ gives
\begin{equation}
 e_{n+1}=(I-\alpha A)e_n,\qquad
 \|e_n\|_A\le
 \left(\frac{\kappa-1}{\kappa+1}\right)^n\|e_0\|_A.
\end{equation}
Using the same synthetic source and coordinate-invariant quadratic-norm
error, tree whitening improves the 16-step relative error from
$1.6710\,10^{-3}$ to $1.3401\,10^{-5}$. A $\pi$ rescaling leaves it
unchanged. The relative-VEV chart worsens it to $.50256$.
This is an actual solver improvement, not a claim of faster physical
loop convergence. The full negative-curvature kernel must not be
passed to an SPD solver just because this local $M$ is positive.

\subsection{Why absorbing Z is not an all-orders repair}
Let $A=G^{-1/2}ZG^{-1/2}$. The familiar expansion
\begin{equation}
 (I+A)^{-1}=\sum_{n=0}^\infty(-A)^n
\end{equation}
fails for the known $A$ because $\rho(A)=2.46659>1$. Directly inverting
this finite matrix is valid as a linear-algebra operation. Declaring
$M=G+Z$ to be a new free metric gives instead
\begin{equation}
 \operatorname{eig}(Z,M)=\frac{\lambda_i}{1+\lambda_i}
       =(.0089150,.2146845,.7115323).
\end{equation}
The apparent improvement is a change of free/interacting split, not
the invariance test above. It sums repeated known insertions but does
not determine independent two-loop self-energies, vertex corrections,
counterterms or the omitted soft nonlocal part. Any proposed resummation
must specify subtraction against double counting and test its next
independent remainder in the new counting. The number $.7115$ is not
an all-orders majorant.

\section{A genuine series improvement from the energy threshold}
The equal-mass one-loop bubble contains the finite momentum difference
\begin{equation}
 F(z)=\int_0^1dx\,\log[1+zx(1-x)],\qquad z=s/m^2,\quad s=p_E^2.
\end{equation}
This is an actual master in the existing kernel. It is dimensionless
and independent of a unit convention for $\pi$ or $\mu$.

\subsection{Derivation of the ordinary series and its obstruction}
For $|z|<4$, expand the logarithm uniformly away from its limiting
singularity and integrate term by term:
\begin{align}
 F(z)&=\sum_{n=1}^\infty
       \frac{(-1)^{n+1}z^n}{n}\int_0^1x^n(1-x)^n\,dx\\
 &=\sum_{n=1}^\infty
       \frac{(-1)^{n+1}(n!)^2}{n(2n+1)!}z^n.
 \label{eq:taylor}
\end{align}
The closest branch point is $z=-4$, where the logarithm's argument
vanishes at $x=1/2$. Thus the Taylor radius is four, determined by
the two-particle threshold, not by a fitted convergence tolerance.
Boundary convergence is a separate question; outside this radius
the ordinary power series cannot converge.

\subsection{Threshold-derived map and exact coefficients}
Map the plane cut along $(-\infty,-4]$ to the unit disk:
\begin{equation}
 w=\frac{\sqrt{1+z/4}-1}{\sqrt{1+z/4}+1},
 \qquad z=\frac{16w}{(1-w)^2}.
 \label{eq:map}
\end{equation}
For every finite Euclidean $z>0$, $0<w<1$.
The branch point fixes the factor four; this is not a free parameter.
A direct evaluation of the integral, using $x=(1+t)/2$, gives
\begin{equation}
 F(z)=2\left[\sqrt{\frac{4+z}{z}}\,
                   \operatorname{arsinh}\frac{\sqrt z}{2}-1\right].
\end{equation}
For positive $w$,
\begin{equation}
 \sqrt{\frac{4+z}{z}}=\frac{1+w}{2\sqrt w},\qquad
 \operatorname{arsinh}\frac{\sqrt z}{2}=2\operatorname{artanh}\sqrt w.
\end{equation}
Substituting $\operatorname{artanh}t/t=\sum_{k\ge0}t^{2k}/(2k+1)$ yields
\begin{equation}
 \boxed{F(z(w))=\sum_{n=1}^\infty\frac{8n}{4n^2-1}w^n.}
 \label{eq:mapped}
\end{equation}
Analytic continuation extends this identity to $|w|<1$. Unlike changing
a ruler, the nonlinear map changes the location of the nearest complex
singularity relative to the series variable.
Threshold-aware conformal representations of correlators have established
precedents \cite{Greynat}; the explicit proof here is for this master.

\begin{proposition}[Certified Euclidean tail]
For $0\le w<1$ and $N\ge1$, set $F_N=\sum_{n=1}^N8n w^n/(4n^2-1)$.
Then
\begin{equation}
 0\le F-F_N\le
 \frac{8(N+1)}{4(N+1)^2-1}\frac{w^{N+1}}{1-w}.
 \label{eq:tail}
\end{equation}
\end{proposition}
\begin{proof}
The coefficients $c_n=8n/(4n^2-1)$ are positive and decrease:
the derivative of $8t/(4t^2-1)$ is
$-8(4t^2+1)/(4t^2-1)^2<0$ for $t\ge1$.
Therefore $\sum_{n>N}c_nw^n\le c_{N+1}\sum_{n>N}w^n$.
For complex $w$, replacing $w$ by $|w|$ gives an absolute bound.
\end{proof}

\subsection{Test with the actual frozen spectrum}
For the lightest massive vector in P-MOM1,
\begin{equation}
 m^2=.00506506436\,\omega^2,\quad s=.05\,\omega^2,\quad
 z=9.87154287,\quad w=.301243161,\quad F=.9243817811.
\end{equation}
The Taylor series is outside its disk; the mapped series is well inside.
The independent integral and exact expression agree. Errors are
absolute errors in this dimensionless master:
\begin{center}
\begin{tabular}{rrrr}
\toprule
Terms & Taylor error & Mapped error & Proven mapped bound\\
\midrule
4 & $2.3811$ & $1.3411\,10^{-3}$ & $1.4345\,10^{-3}$\\
8 & $34.6267$ & $6.2549\,10^{-6}$ & $6.5172\,10^{-6}$\\
16 & $1.7732\,10^4$ & $2.2811\,10^{-10}$ & $2.3347\,10^{-10}$\\
32 & $1.2175\,10^{10}$ & $5.4594\,10^{-19}$ & $5.5279\,10^{-19}$\\
\bottomrule
\end{tabular}
\end{center}
No coefficient was fitted to these values. The script also tests all six
existing momentum probes and the certified P50 mass at unchanged $\tau=1$.
The table measures series truncation at the fixed stored input, not
uncertainty in that mass or the accuracy of a physical prediction.
It does not generate another quartic or scalar-background scan.
P-MOM1 already evaluates the exact masters; its values do not need this
series to be repaired. The benefit is a controlled alternative expansion,
compression and potentially analytic-continuation strategy.

\section{Other structural options and the decision boundary}
\paragraph{Retain near-threshold modes rather than a local Taylor series.}
For unequal masses the normal two-particle threshold is
$s_{\rm th}=(\sqrt a+\sqrt b)^2$, with $a,b$ squared masses.
A useful expansion test is $|s|/s_{\rm th}$ rather than the numerical
size of a mass in arbitrary units. Retain modes near this threshold, or
their full nonlocal kernels, when the local derivative expansion is
not controlled. The separation of hard and soft contributions must
remain consistent with matching \cite{Dittmaier}.
At the current certified transverse ray,
\begin{equation}
 4m_{50}^2=48\tau\sigma^2,\qquad
 4m_{36}^2=\frac{232}{5}\tau\sigma^2.
\end{equation}
Weakening transverse quartics therefore shrinks the momentum Taylor disk
at the same time that the certified vertex-to-gap loop bound grows.
Changing units cannot fix either effect. This strengthens the reason
not to use another blind quartic reduction.

\paragraph{Keep exact massless nonanalytic pieces.}
A doubly massless scalar bubble contains $\log(s/\mu^2)-2$ and has no
analytic Taylor expansion about $s=0$. Do not apply \eqref{eq:map}
with $m=0$. Separate the known log exactly and approximate only the
remaining analytic part. Tensor reductions may contain cancellations
and gauge-artifact thresholds; approximations must preserve common
counterterms and Ward/Nielsen identities rather than approximate each
large term independently without an assembled error bound.

\paragraph{Distinguish a momentum map from Borel resummation.}
The proof above concerns powers of momentum at a fixed loop order,
not powers of a quantum coupling. Borel-plane conformal resummation
addresses a different question and needs large-order/singularity
information. Recent QCD work \cite{Caprini} cannot be imported as a
convergence theorem for this P54 action. One known loop matrix does not
determine a Borel transform or a controlled all-orders sum.

\paragraph{Recommended order.}
Use tree-metric whitening as a numerical convention while always
reporting invariant generalized spectra. If a series representation is
needed, extend the certified threshold map to unequal-mass and mixed
masters, retaining massless logs exactly and certifying the assembled
kernel error. This is a bounded auxiliary task, not a substitute for
the fixed-bare Nielsen/BRST and complete physical mixing analysis.
Retained-light-block EFT power counting is the structural alternative.
Physical fit, full matching and finite seesaw gates stay unchanged.

\section{Reproducibility}
The verifier \path{route_f/code/verify_p54_coordinate_convergence.py}
reads the saved P-MOM1 and common-renormalization JSONs, checks the former's
source hashes, and uses NumPy/SciPy and 60-digit mpmath. It passes
183/183 coordinate, source-response, solver and spectral tests.
No new action Hessian is evaluated. Its JSON distinguishes proved
matrix/series identities from numerical evidence and explicitly sets
all-orders loop convergence and full-kernel conformal replacement to false.

\clearpage
\begin{thebibliography}{9}
\bibitem{Criado}
J.~C.~Criado and M.~Perez-Victoria,
\emph{Field redefinitions in effective theories at higher orders},
\href{https://arxiv.org/abs/1811.09413}{arXiv:1811.09413}.
\bibitem{Greynat}
D.~Greynat and S.~Peris,
\emph{Resummation of Threshold, Low- and High-Energy Expansions for
Heavy-Quark Correlators},
\href{https://arxiv.org/abs/1006.0643}{arXiv:1006.0643}.
\bibitem{Dittmaier}
S.~Dittmaier, S.~Schuhmacher and M.~Stahlhofen,
\emph{Integrating out heavy fields in the path integral using the
background-field method: general formalism},
\href{https://arxiv.org/abs/2102.12020}{arXiv:2102.12020}.
\bibitem{Caprini}
I.~Caprini, \emph{Revisiting the convergence of the perturbative QCD
expansions based on conformal mapping of the Borel plane},
\href{https://arxiv.org/abs/2312.07143}{arXiv:2312.07143} (2023).
\end{thebibliography}
\end{document}
